% !TeX spellcheck = he_IL
%% ! TEX TS-program = xelatexmk
% above: commands appropriate for TeXShop in mac OS
\documentclass[12pt, a4paper]{article}
\usepackage{etoolbox}
\usepackage{bbm}
%\usepackage{pifont}

%%%%%%%%%%%%%%%%%%%%%%%%%%
%%
%%
%%  Exam Template
%%
%%
%%%%%%%%%%%%%%%%%%%%%%%%%%
\newtoggle{testreview}
\toggletrue{testreview}  
\togglefalse{testreview} % uncomment to make it ready for exam unit
\newtoggle{testsolution}
\toggletrue{testsolution}  
\togglefalse{testsolution} % comment to show solutions
\newcounter{version}
\setcounter{version}{2}  % version of multiple-choice

	







\usepackage[a4paper, top=3cm,left=2.5cm, right=2.3cm, bottom=2.7cm]{geometry}
\usepackage{amsmath, amssymb, array, multicol, enumitem}
\usepackage{pdfpages}
\usepackage{comment, graphicx}  
\usepackage{xparse}
\DeclareGraphicsExtensions{.png,.jpg}
\usepackage{pgffor}
\usepackage[colorlinks=true]{hyperref} 
\AtBeginDocument{%
	\hypersetup{pdfborder={0 0 0},
		urlcolor=blue}%
}

\usepackage{wasysym}
\usepackage{fontspec}
\usepackage{xunicode}
\usepackage{polyglossia}
\setdefaultlanguage{hebrew}
\setotherlanguage{english}
%\newfontfamily\hebrewfont[Script=Hebrew, Ligatures=TeX, Scale=1.1, WordSpace =1.05]{David}
%\newfontfamily\hebrewfont[Script=Hebrew, Ligatures=TeX, Scale=1.1, WordSpace=1.05]{David  CLM}
\newfontfamily\hebrewfont[Script=Hebrew]{Times New Roman}

%http://tex.stackexchange.com/questions/141231/equation-number-enclosing-parentheses-misplaced-with-polyglossiahebrewamsmath
\makeatletter
%\def\maketag@@@#1{\hbox{\m@th\normalfont\LRE{#1}}}
\def\tagform@#1{\maketag@@@{(\ignorespaces#1\unskip)}}
\makeatother

%
% Hebrew fonts are two small 
%
\fontsize{14}{15}

%
% Package fro proper positioning of logo
%
\usepackage[absolute]{textpos}


% http://tex.stackexchange.com/questions/10589/fix-nested-section-numbers-in-rtl-languages-with-polyglossia/
\makeatletter
\renewcommand{\SepMark}[1]{\def\@SepMark{#1}}
\makeatother
\SepMark{.}
\robustify\hebrewnumeral
\robustify\Hebrewnumeral
\robustify\Hebrewnumeralfinal


% from polyglossia docs
%\gappto\captionshebrew{%
\def\urladdrname{{\itshape כתובת רשת}}
\def\emailaddrname{{\itshape כתובת דואל}}
%}

% https://tex.stackexchange.com/questions/176661/polyglossia-arabic-text-english-bibliography
\makeatletter
\let\orig@MakeUppercase \MakeUppercase
\makeatother
\providecommand{\MakeUppercase}[1]{#1}




%%%%%%%%%%%%%%%%%%%%%%%%%%
%%
%%
%%   User Commands
%%
%%
%%%%%%%%%%%%%%%%%%%%%%%%%%
\newcommand{\heb}[1]{\begin{small}\R{#1}\end{small}}
\newcommand{\hebt}[1]{\begin{scriptsize}\R{#1}\end{scriptsize}}
\newcommand{\cc}{\mathbb{C}}
\newcommand{\nn}{\mathbb{N}}
\newcommand{\zz}{\mathbb{Z}}
%\renewcommand{\phi}{\varphi}
\newcommand{\abs}[1]{\left|#1\right|}
\newcommand{\lr}[1]{\left(#1\right)}


\newcommand{\ff}{\mathbb{F}}
\newcommand{\FF}{\mathbb{F}}
\newcommand{\mtn}{m\times n}
\newcommand{\longr}{\longrightarrow}
\newcommand{\longrl}{\longleftrightarrow}
\newcommand{\ol}{\overline}
\newcommand{\ntn}{n\times n}
\newcommand{\rr}{\mathbb{R}}
\newcommand{\QQ}{\mathbb{Q}}
\newcommand{\RR}{\mathbb{R}}
\newcommand{\im}{\mathrm{Im}}




\newcommand{\twopartdef}[4]
{
	\left\{
		\begin{array}{ll}
			#1 & \mbox{   } #2 \\
			#3 & \mbox{   } #4
		\end{array}
	\right.
}

\newcommand{\solutionword}{\textbf{פתרון:}}
\newcommand{\points}[1]{(#1 נק') }
\iftoggle{testsolution}{%
	\newcommand{\sols}[1]{\solutionword$\;\;$ #1$\square$}
	\newcommand{\spacing}{0cm}
	\newcommand{\questionspacing}{2.5cm}
	\newcommand{\extrapages}{}
}{ %
%\newcommand{\spacing}{12cm}
%\newcommand{\halfspacing}{0.5cm}	
\newcommand{\spacing}{24cm}
\newcommand{\questionspacing}{\spacing}
\newcommand{\sols}[1]{}
\newcommand{\extrapages}{%
	\newpage
	\noindent
	דף זה ריק במכוון ומיועד לשימוש הסטודנט/ית.
	\newpage
	\noindent
	דף זה ריק במכוון ומיועד לשימוש הסטודנט/ית.
}
\iftoggle{testreview}{%
\renewcommand{\spacing}{1cm}
%\newcommand{\questionspacing}{\spacing}
\renewcommand{\extrapages}{}	
}{}	
}






%%%%%%%%%%%%%%%%%%%%%%%%
%%
%%   BEGIN DOCUMENT
%%
%%%%%%%%%%%%%%%%%%%%%%%%
\usepackage{graphicx}

\begin{document}

\selectlanguage{hebrew}


\newcommand{\semester}{%
	סמסטר א' תשפ״ו
}
\newcommand{\course}{%
 אלגברה לינארית להנדסה 
}
\newcommand{\coursenumber}{9321}	
\newcommand{\lecturers}{%
	מרצים: ד״ר י.מייזל, ד״ר נ.גולקו, ד״ר נ.מאיר, ד״ר נ.איידלשטין, ד״ר מ.פורת, ד״ר מ.זיו-אב, ד״ר מ.ברשדסקי
}

\newcommand{\testdate}{{12.12.25}}
\iftoggle{testsolution}{%
	\newcommand{\testtype}{%
		פתרונות למבחן מועד א'
	}
}{%
	\newcommand{\testtype}{%
	\boxed{\text{
	בוחן לדוגמה
	}	}
}}

\newcommand{\duration}{%
	2 שעות.
}
\newcommand{\helpmaterial}{%
מחשבון פשוט בלבד.
}



\thispagestyle{empty}

\begin{textblock*}{9cm}[0.5,0.5](90mm,3cm)
	\includegraphics[width=2in]{logo}
\end{textblock*}	

\noindent

\begin{center}
אוניברסיטת בן-גוריון בנגב - היחידה להוראת קורסי יסוד - %
\smallskip
\semester \\
{\large \bf \course - 214.1.\coursenumber}\\
\smallskip
\lecturers\\
\smallskip
\smallskip
\smallskip
{\Large \testtype }\\
\end{center}

\noindent
\textbf{%
תאריך:
} \testdate\\
\textbf{משך הבחינה: 
} \duration \\
\textbf{
חומר עזר: 
}
\helpmaterial



\noindent

\begin{itemize}
\item
\textbf{
בבוחן 10 שאלות רב-ברירה. יש לענות על כולן. כל שאלה במשקל של 10 נקודות.}
\item
לכל שאלה יש רק תשובה אחת נכונה בדיוק. 
\item
יש להשחיר תשובה נכונה אחת בלבד לכל שאלה בטבלה. 
\item
בכל שאלה בה תסומן יותר מתשובה אחת או שלא תסומן אף תשובה, הניקוד עבור השאלה יהיה
$0$.
\iffalse
\item התשובות שלי הן:  

\begin{center}
\begin{tabular}{|c|c|c|c|c|}
\hline
1 & 2 &  3 & 4 & 5 \\
\hline
\begin{tabular}{cccc}
א & ב & ג &ד
\end{tabular}
&
\begin{tabular}{cccc}
א & ב & ג &ד
\end{tabular}

&
\begin{tabular}{cccc}
א & ב & ג &ד

\end{tabular}
&
\begin{tabular}{cccc}
א & ב & ג &ד
\end{tabular}
&
\begin{tabular}{cccc}
א & ב & ג &ד
\end{tabular}
\\
\hline
\end{tabular}


\begin{tabular}{|c|c|c|c|c|}
\hline
 6 & 7& 8& 9& 10  \\
\hline
\begin{tabular}{cccc}
א & ב & ג &ד
\end{tabular}
&
\begin{tabular}{cccc}
א & ב & ג &ד
\end{tabular}
&
\begin{tabular}{cccc}
א & ב & ג &ד
\end{tabular}
&
\begin{tabular}{cccc}
א & ב & ג &ד
\end{tabular}
&
\begin{tabular}{cccc}
א & ב & ג &ד
\end{tabular}
\\
\hline
\end{tabular}
\end{center}


\item התשובות שלי הן:  

\begin{center}
\begin{tabular}{|c|c|c|c|c|c|c|c|c|c|}
\hline
1 & 2 &  3 & 4 & 5 & 6 & 7 & 8 & 9 & 10 \\
\hline
\begin{tabular}{cccc}
א\\
ב\\
ג\\
ד\\
\end{tabular}
&
\begin{tabular}{cccc}
א\\
ב\\
ג\\
ד\\
\end{tabular}

&
\begin{tabular}{cccc}
א\\
ב\\
ג\\
ד\\

\end{tabular}
&
\begin{tabular}{cccc}
א\\
ב\\
ג\\
ד\\
\end{tabular}
&
\begin{tabular}{cccc}
א\\
ב\\
ג\\
ד\\
\end{tabular}
&
\begin{tabular}{cccc}
א\\
ב\\
ג\\
ד\\
\end{tabular}
&
\begin{tabular}{cccc}
א\\
ב\\
ג\\
ד\\
\end{tabular}
&
\begin{tabular}{cccc}
א\\
ב\\
ג\\
ד\\
\end{tabular}
&
\begin{tabular}{cccc}
א\\
ב\\
ג\\
ד\\
\end{tabular}
&
\begin{tabular}{cccc}
א\\
ב\\
ג\\
ד\\
\end{tabular}
\\
\hline
\end{tabular}
\end{center}


\fi


\item
\textbf{%
שימו לב:
 דפי הטיוטה נשלחים לגריסה באופן מיידי. 
}  

\noindent\rule{\textwidth-30pt}{0.4pt}


\item
$M_{m\times n}(\FF)$
זו קבוצת המטריצות בגודל
$m\times n$
(
$m$
שורות,
$n$
עמודות
)
עם רכיבים מ-$\ff$.
%%%%%



\item
לכל מטריצה
$A$,
נסמן
על ידי
$(A)_{i,j}$
את האיבר של 
$A$
הנמצא בשורה ה-$i$
ובעמודה ה-$j$.
%%%%%



\item
 פיתוח דטרמיננטה של $A\in M_{n\times n}(\FF)$ לפי השורה ה-$i$ 
או לפי העמודה ה-$j$:
$$\det(A)=\sum_{i=1}^n (-1)^{i+j}a_{ij}\det(M_{ij}(A))=\sum_{j=1}^n (-1)^{i+j}a_{ij}\det(M_{ij}(A))$$
כאן $M_{ij}(A)$ הוא המינור ה-$i,j$ של המטריצה $A$. 


\noindent\rule{\textwidth-30pt}{0.4pt}

\begin{center}
{\Large
	בהצלחה!
}
\end{center}

\begin{center}
\begin{tabular}{|c||c||c||c||c||c||c||c||c||c|}
\hline
1 & 2 &  3 & 4 & 5 & 6 & 7 & 8 & 9 & 10 \\
\hline
\begin{tabular}{cccc}
\\\textcircled{א}\\\\
\textcircled{ב}\\\\
\textcircled{ג}\\\\
\textcircled{ד}\\\\
\end{tabular}
&
\begin{tabular}{cccc}
\\\textcircled{א}\\\\
\textcircled{ב}\\\\
\textcircled{ג}\\\\
\textcircled{ד}\\\\
\end{tabular}

&
\begin{tabular}{cccc}
\\\textcircled{א}\\\\
\textcircled{ב}\\\\
\textcircled{ג}\\\\
\textcircled{ד}\\\\
\end{tabular}
&
\begin{tabular}{cccc}
\\\textcircled{א}\\\\
\textcircled{ב}\\\\
\textcircled{ג}\\\\
\textcircled{ד}\\\\
\end{tabular}
&
\begin{tabular}{cccc}
\\\textcircled{א}\\\\
\textcircled{ב}\\\\
\textcircled{ג}\\\\
\textcircled{ד}\\\\
\end{tabular}
&
\begin{tabular}{cccc}
\\\textcircled{א}\\\\
\textcircled{ב}\\\\
\textcircled{ג}\\\\
\textcircled{ד}\\\\
\end{tabular}
&
\begin{tabular}{cccc}
\\\textcircled{א}\\\\
\textcircled{ב}\\\\
\textcircled{ג}\\\\
\textcircled{ד}\\\\
\end{tabular}
&
\begin{tabular}{cccc}
\\\textcircled{א}\\\\
\textcircled{ב}\\\\
\textcircled{ג}\\\\
\textcircled{ד}\\\\
\end{tabular}
&
\begin{tabular}{cccc}
\\\textcircled{א}\\\\
\textcircled{ב}\\\\
\textcircled{ג}\\\\
\textcircled{ד}\\\\
\end{tabular}
&
\begin{tabular}{cccc}
\\\textcircled{א}\\\\
\textcircled{ב}\\\\
\textcircled{ג}\\\\
\textcircled{ד}\\\\
\end{tabular}
\\
\hline
\end{tabular}
\ 

\ 

\huge{
הוראות למילוי הטבלה בעמוד הבא
}
\end{center}
\pagebreak

\section*{הוראות למילוי הטבלה}
\begin{itemize}
    \item[$\bullet$] יש להשתמש בעט שחור או כחול. 
    \item[$\bullet$] יש לסמן תשובה אחת בלבד ע"י השחרת העיגול המתאים.
    \item[$\bullet$]
    בשאלה בה לא יושחר אף עיגול,
    הניקוד יהיה 0.
    \item[$\bullet$] בשאלה בה יושחרו יותר מעיגול אחד, הניקוד יהיה 0.
.
    \item[$\bullet$] כדי לתקן סימון, יש לסמן X
    על התשובה הלא נכונה ולהשחיר את התשובה הנכונה.

    \item[$\bullet$]
    \textbf{שימו לב:}
    לא ניתן לבחור סעיף שסומן עליו X
    לכן הקפידו לסמן את תשובתכם כשאתם בטוחים בה.
\end{itemize}

\section*{דוגמאות}


\begin{center}
\begin{tabular}{|c||c||c||c||c||c||c||c||c||c|}
\hline
1 & 2 &  3 & 4 & 5 & 6 & 7 & 8 & 9 & 10 \\
\hline
\begin{tabular}{c}
\\\Large{\CIRCLE}\\\\
\ \  \hspace{-8pt}\raisebox{0pt}[0pt][0pt]{\Large\CIRCLE}%
\  \hspace{4pt}\llap{\raisebox{-4pt}[0pt][0pt]{ $\scalebox{3}{$\times$}$}}\hspace{-8pt}
\\\\
\textcircled{ג}\\\\
\textcircled{ד}\\\\
\end{tabular}
&
\begin{tabular}{cccc}
\\\textcircled{א}\\\\
\Large{\CIRCLE}\\\\
\textcircled{ג}\\\\
\textcircled{ד}\\\\
\end{tabular}
&
\begin{tabular}{cccc}
\\
\textcircled{א}
\\\\
\ \  \hspace{-8pt}\raisebox{0pt}[0pt][0pt]{\Large\CIRCLE}%
 \ \hspace{4pt}\llap{\raisebox{-4pt}[0pt][0pt]{ $\scalebox{3}{$\times$}$}}\hspace{-8pt}
\\\\
\Large{\CIRCLE}\\\\
\ \  \hspace{-8pt}\raisebox{0pt}[0pt][0pt]{\Large\CIRCLE}%
 \ \hspace{4pt}\llap{\raisebox{-4pt}[0pt][0pt]{ $\scalebox{3}{$\times$}$}}\hspace{-8pt}
\\\\
\end{tabular}
&
\begin{tabular}{cccc}
\\\textcircled{א}\\\\
\textcircled{ב}\\\\
\textcircled{ג}\\\\
\textcircled{ד}\\\\
\end{tabular}
&
\begin{tabular}{cccc}
\\\hspace{-8pt}\ \  \raisebox{0pt}[0pt][0pt]{\Large\CIRCLE}%
 \ \hspace{4pt}\llap{\raisebox{-4pt}[0pt][0pt]{ $\scalebox{3}{$\times$}$}}\hspace{-8pt}
\\\\
\textcircled{ב}\\\\
\textcircled{ג}\\\\
\textcircled{ד}\\\\
\end{tabular}
&
\begin{tabular}{cccc}
\\\hspace{-8pt}\ \  \raisebox{0pt}[0pt][0pt]{\Large\CIRCLE}%
 \ \hspace{4pt}\llap{\raisebox{-4pt}[0pt][0pt]{ $\scalebox{3}{$\times$}$}}\hspace{-8pt}
\\\\
\hspace{-8pt}\ \  \raisebox{0pt}[0pt][0pt]{\Large\CIRCLE}%
 \ \hspace{4pt}\llap{\raisebox{-4pt}[0pt][0pt]{ $\scalebox{3}{$\times$}$}}\hspace{-8pt}
\\\\
\textcircled{ג}\\\\
\textcircled{ד}\\\\
\end{tabular}
&
\begin{tabular}{cccc}
\\\Large{\CIRCLE}\\\\
\Large{\CIRCLE}\\\\
\textcircled{ג}\\\\
\textcircled{ד}\\\\
\end{tabular}
&
\begin{tabular}{cccc}
\\\Large{\CIRCLE}\\\\
\Large{\CIRCLE}\\\\
\Large{\CIRCLE}\\\\
\textcircled{ד}\\\\
\end{tabular}
&
\begin{tabular}{cccc}
\\\hspace{-8pt}\ \  \raisebox{0pt}[0pt][0pt]{\Large\CIRCLE}%
 \ \hspace{4pt}\llap{\raisebox{-4pt}[0pt][0pt]{ $\scalebox{3}{$\times$}$}}\hspace{-8pt}
\\\\
\Large{\CIRCLE}\\\\
\textcircled{ג}\\\\
\Large{\CIRCLE}\\\\
\end{tabular}
&
\begin{tabular}{cccc}
\\\hspace{-8pt}\ \  \raisebox{0pt}[0pt][0pt]{\Large\CIRCLE}%
 \ \hspace{4pt}\llap{\raisebox{-4pt}[0pt][0pt]{ $\scalebox{3}{$\times$}$}}\hspace{-8pt}
\\\\
\Large{\CIRCLE}
\\\\
\hspace{-8pt}\ \  \raisebox{0pt}[0pt][0pt]{\Large\CIRCLE}%
 \ \hspace{4pt}\llap{\raisebox{-4pt}[0pt][0pt]{ $\scalebox{3}{$\times$}$}}\hspace{-8pt}
\\\\
\Large{\CIRCLE}\\\\
\end{tabular}
\\
\hline
\end{tabular}
\ 

\ 

\begin{itemize}[label=$\bullet$]
    \item שאלות 1,2,3 עם סימונים תקינים עבור תשובות:
    \begin{enumerate}
        \item א
        \item ב
        \item ג
    \end{enumerate}
    \item בשאלות 4,5,6 חסר סימון והניקוד יהיה 0.
    \item בשאלות 7,8,9,10
    יש יותר מסימון אחד והניקוד יהיה 0
\end{itemize}


\end{center}

\iffalse
\item
\textbf{\underline{תזכורת:}}
$$,I_p=
\left\{
f:\rr\to\cc
\mid
\text{$f$ מחזורית-$2\pi$ ואינטגרבילית}
\right\}$$
$$,PC_p=
\left\{
f:\rr\to\cc
\mid
\text{$f$ מחזורית-$2\pi$ ורציפה למקוטעין}
\right\}$$
$$.C_p=
\left\{
f:\rr\to\cc
\mid
\text{$f$ מחזורית-$2\pi$ ורציפה}
\right\}$$

\item
\textbf{\underline{תזכורת:}}
סדרת פונקציות ממשיות
$(F_n)_{n=1}^{\infty}$
ב-$PC_p$
נקראת גרעין טוב (סומוביליות חיובי) אם 
\begin{itemize}
\item
לכל
$n\ge1$
מתקיים
$F_n(x)\ge0$
לכל
$x\in[-\pi,\pi]$.
\item
לכל
$n\ge1$
מתקיים
$\frac{1}{2\pi}
\int_{-\pi}^{\pi}
F_n(x)dx=1$.
\item
לכל
$0<\delta<\pi$
מתקיים
$\underset{n\to\infty}{\lim}
\frac{1}{2\pi}
\int_{-\delta}^{\delta}
F_n(x)dx=1$.
\end{itemize}

\item
\textbf{\underline{תזכורת:}}
לכל
$N\ge0$ 
ולכל
$f\in I_p$:
$$.\left(S_N(f)\right)(x)
=\sum_{n=-N}^N \widehat{f}(n)e^{inx}
\quad,\quad
\left(\sigma_N(f)\right)(x)=
\sum_{n=-N}^N
\left(
1-\frac{|n|}{N+1}\right)\widehat{f}(n)e^{inx}$$

\item
\textbf{\underline{תזכורת:}}
עבור קטע
$[c,d]\subset\rr$,
נסמן את הפונקציה המציינת שלו על ידי
$$.\mathbbm{1}_{[c,d]}(x)=
\begin{cases}
1&: x\in[c,d]\\
0&: \text{אחרת}
\end{cases}$$
\noindent\rule{\textwidth-30pt}{0.4pt}

\fi
\end{itemize}

%%%
%%%
%%%   End of Header and Instructions
%%%
%%%

\newpage

\underline{{\Large
}}


\begin{enumerate}

%%%%%%%%%%%%%%%%%%%%%%%%%%%%%%%%%%%%%%
%%  Problems
%%%%%%%%%%%%%%%%%%%%%%%%%%%%%%%%%%%%%%


%%%%% שאלה 2 %%%%%%%


\item\points{10}
נתון ישר ב-$\rr^2$
שעובר דרך הנקודות
$(1,2)$
ו-
$(3,5)$.
מי מהבאות היא הצגה פרמטרית של הישר?
\begin{enumerate}
\item
$L=\{ (-1,-2)+t\cdot (-4,-6)\bigm|
t\in\rr\}$.
\item
$L=\{ (-3,-4)+t\cdot (4,-6)\bigm|
t\in\rr\}$.
\item
$L=\{ (-3,-4)+t\cdot (4,6)\bigm|
t\in\rr\}$.
\item
$L=\{ (-1,-2)+t\cdot (-4,6)\bigm|
t\in\rr\}$.
\end{enumerate}

\noindent\rule{\textwidth-30pt}{0.4pt}

\item\points{10}
נתון מישור
$H$
במרחב
$\rr^3$
העובר דרך הנקודה
$(2,-1,1)$,
שמקביל למישור הנתון על ידי המשוואה
$x+2y+z=3$.
בחרו את הטענה הנכונה:
\begin{enumerate}
\begin{multicols}{2}
\item
$(0,0,0)\in H$
\item
$(3,-2,2)\in H$
\item
$(5,0,-6)\in H$
\item
$(0,1,0)\in H$
\end{multicols}
\end{enumerate}

\noindent\rule{\textwidth-30pt}{0.4pt}

%%%%% שאלה 2 %%%%%%%

\item\points{10}
תהי
$A\in M_{m\times n}(\rr)$
ויהי
$b\in \rr^m$.
אילו מהטענות הבאות בהכרח נכונה?
\begin{enumerate}
\item
אם
$m>n$
אז לממ״ל
$AX=b$
אין פתרון.
%%%
\item
אם
$n>m$
אז לממ״ל
$AX=0$
יש פתרון לא טריוויאלי.
%%%
\item
אם קיים פתרון לממ״ל
$AX=b$
אז לכל
$c\in \rr^m$
קיים פתרון לממ״ל
$AX=c$.
%%%
\item
אם
$u,v$
הם פתרונות של הממ״ל
$AX=b$
אז גם הוקטור
$2u+v$
הוא פתרון של אותה הממ״ל.
\end{enumerate}

\noindent\rule{\textwidth-30pt}{0.4pt}


\item\points{10}
נתונה ממ״ל התלויה בפרמטר ממשי
$a$:
$$\begin{cases}
x_1\,\,\,+\quad \quad\,\,\, \, ax_2\,+\,\,\,\,\,(a+1)x_3&=a\\
ax_1+(a^2+1)x_2+(a^2+2a)x_3&=a^2-1\\
\quad\quad \quad\quad\quad\quad x_2+\quad\quad\quad\,\, \,2x_3&=a\\
\end{cases}$$
אילו מהטענות הבאות בהכרח נכונה?
\begin{enumerate}
\begin{multicols}{2}
\item
אם קיים פתרון לממ״ל אז הוא יחיד.
\item
לכל
$a\in\rr$
קיים לממ״ל פתרון.
\item
אם
$a=-1$
אז לממ״ל יש אינסוף פתרונות.
\item
לכל
$a\in\rr$
לממ״ל אין פתרון.
\end{multicols}
\end{enumerate}

\noindent\rule{\textwidth-30pt}{0.4pt}


%%%
%%%%%%
%%%

\item\points{10}
תהי
$A\in M_{4\times 3}(\cc)$
מטריצה מדורגת קנונית, עבורה
$\mathrm{rank}(A)=2$
כך ש-
$(1,0,-2)$
פתרון של הממ״ל
$AX=0$.
למה שווה סכום כל האיברים בעמודה השלישית של
$A$:
\begin{enumerate}
\begin{multicols}{4}
\item
$1$
\item
$-2$
\item
$0$
\item
$\frac{1}{2}$
\end{multicols}
\end{enumerate}

\newpage


%%%
%%%%%%
%%%



\item\points{10}
נתונות המטריצות הבאות
$$A=\begin{pmatrix}
1&2+i\\i+1&2\\
\end{pmatrix}\quad,\quad
B=\begin{pmatrix}
2i&1\\0&i\\
\end{pmatrix}$$
ותהי
$C\in M_{2\times 2}(\cc)$
מטריצה  המקיימת
$A\cdot C=B$.
בחרו את הטענה הנכונה:
\begin{enumerate}
\begin{multicols}{2}
\item
$(C)_{1,1}-(C)_{2,1}=2$
\item
$(C)_{1,1}-(C)_{2,1}=-2$
\item
$(C)_{1,1}-(C)_{2,1}=i$
\item
$(C)_{1,1}-(C)_{2,1}=-i$
\end{multicols}
\end{enumerate}

\noindent\rule{\textwidth-30pt}{0.4pt}


%%%
%%%%%%
%%%

\item\points{10}
מה הטענות הבאות בהכרח נכונה?
\begin{enumerate}
\item
אם
$A\in M_{3\times 3}(\rr)$
אלכסונית, אז לכל
$B\in M_{3\times 3}(\rr)$
מתקיים
$AB=BA$.
\item
אם
$E_1,E_2\in M_{3\times 3}(\rr)$
מטריצות אלמנטריות, אז
$E_1\cdot E_2= E_2\cdot E_1$.
\iffalse
\item
אם
$A\cdot B$
משולשית עליונה, אז
$A$
משולשית עליונה וגם
$B$
משולשית עליונה.
\fi
\item
אם 
$E\in M_{3\times3}(\rr)$
אלמנטרית,
אז
$E^t$
היא אלמנטרית.
\item
אם
$A\in M_{3\times 3}(\rr)$
מקיימת
$A^2=0_{3\times 3}$,
אז
$A=0_{3\times 3}$.
\end{enumerate}


%%%
%%%%%%
%%%

\noindent\rule{\textwidth-30pt}{0.4pt}

\item\points{10}
תהיינה
$A,B\in M_{\ntn}(\rr)$
מטריצות המקיימות
$\det(A)=-1$,
$\det(B)=2$
וגם

$\det(2B^{-1}AB^tA^2B)=-16$.
בחרו את הטענה הנכונה:
\begin{enumerate}
\begin{multicols}{4}
\item
$n=1$
\item
$n=2$
\item
$n=3$
\item
$n=4$
\end{multicols}
\end{enumerate}


\noindent\rule{\textwidth-30pt}{0.4pt}


%%%
%%%%%%
%%%

\item\points{10}
תהי
$$.A=\begin{pmatrix}
1&2&1\\
1&3&1\\
3&7&4\\
\end{pmatrix}$$
ונסמן
$.c= (A^{-1})_{1,1}+(A^{-1})_{2,2}+(A^{-1})_{3,3}$
\iffalse
בעזרת מציאת
$A^{-1}$
(או בכל דרך אחרת בה תבחרו), 
\fi
בחרו את הטענה הנכונה:
\begin{enumerate}
\begin{multicols}{4}
\item
$c=1$
\item
$c=3$
\item
$c=5$
\item
$c=7$
\end{multicols}
\end{enumerate}

\noindent\rule{\textwidth-30pt}{0.4pt}


%%%
%%%%%%
%%%

%%%
%%%%%%
%%%

\item\points{10}
תהיינה
$$A=\begin{pmatrix}
1&2&3\\
0&1&4\\
0&0&1\\
\end{pmatrix}\quad,
B=
\begin{pmatrix}
3&2&1\\
4&1&0\\
-4&0&0\\
\end{pmatrix},\quad
C=
\begin{pmatrix}
1&0&2\\
-2&1&0\\
2&0&4\\
\end{pmatrix}$$
בחרו את הטענה הנכונה:
\begin{enumerate}
\begin{multicols}{2}
\item
$\det(ABC)=39$
\item
$\det(A+B+C)=39$
\item
$\det(CBA)=39$
\item
$\det(B^tAC^t)=39$
\end{multicols}
\end{enumerate}

\iffalse
\newpage

\item\points{10}
מצאו את המטריצה
$A\in M_{3\times 3}(\rr)$
המקיימת את התנאים הבאים:
\begin{itemize}
\item
$A$
אנטי-סימטרית,
\item
הוקטור
$(1,1,1)$
הוא פתרון של הממ״ל ההומוגנית
$AX=0$.
\iffalse
\item
סכום העמודות של
$A$
שווה ל-$\begin{pmatrix}0\\0\\0\\
\end{pmatrix}$,
\fi
\item
$(A)_{1,3}=1$.
\end{itemize}

תשובה: 
$$.A=$$

\noindent\rule{\textwidth-30pt}{0.4pt}

\fi
%%%
%%%%%%
%%%

\underline{
הגדרה:
}
מטריצה $A\in M_{2\times 2}(\rr)$
נקראת סבירה אם מתקיימים התנאים הבאים:
\begin{itemize}
    \item $A$ מדורגת קנונית.
    \item $A\neq O_{2\times 2}$.
    \item $A^2 = O_{2\times 2}$.
\end{itemize}
מה מהסעיפים הבאים נכון?
\begin{enumerate}
    \item לא קיימת מטריצה סבירה.
    \item קיימת מטריצה סבירה אחת ויחידה.
    \item קיימות בדיוק שתי מטריצות סבירות.
    \item קיימות אינסוף מטריצות סבירות.
\end{enumerate}


נתונות מטריצות מדורגות קנונית
$A,B\in M_{2\times 2}(\rr)$ כך ש- $A+B=I_2$.

מה מהבאים 
\underline{בהכרח}
נכון?
\begin{enumerate}
\begin{multicols}{4}
    \item $AB=A$
    \item $BA=B$
    \item $AB=I_2$
    \item $BA=O_{2\times 2}$
\end{multicols}
\end{enumerate}

\end{enumerate}


\end{document}